\section{Overview of the proofs}\label{sec:overview}
We sketch the proofs of Theorems~\ref{thm:bounded} and~\ref{thm:forests} in Section~\ref{sec:sketch-extract}, and that of Theorem~\ref{thm:gap} in Section~\ref{sec:sketch-exact}.

\subsection{Removing rainbow forests one at a time}\label{sec:sketch-extract}
Let $G$ be a graph with $km$ edges and a globally $m$-bounded colouring. We remove rainbow forests with $k$ edges, which we call \emph{$k$-forests}, one at a time.
When $q$ forests are still to be found, the remaining edge set $R$ has $qk$ edges and at most $q$ edges of each colour.
At each step, we choose a random rainbow $k$-forest $T\subseteq R$ in which every edge of $R$ appears with probability (almost exactly) $1/q$, and every colour with $q$ edges in $R$ appears.
By Lemma~\ref{lem:edmonds}, such a $T$ exists whenever $e_R(X)\le q(|X|-1)$ for every nonempty vertex set $X$: then $\one_R/q$ is a convex combination of rainbow $k$-forests.
The difficulty is to keep these conditions true as edges are removed, and here the Laplacian is useful.
If $L_R/q$ is close to a fixed matrix $B$, then, by~\eqref{eq:edgecount}, $e_R(X)/q$ is close to $(\sum_{v\in X}B_{vv}-\one_X^{\!\top}B\one_X)/2$ for every $X$, so a single spectral condition controls all the conditions $e_R(X)\le q(|X|-1)$ at once.
Since every edge is removed with probability $1/q$, we have $\E L_{R\setminus T}/(q-1)=L_R/q$, so $L_R/q$ is a martingale.
Its increments have norm about $\|L_T\|/q$, and, since $L_T\succeq0$, its variance is at most $\|L_T\|\,\E L_T/q^2$.
Thus, once we know the probability that each edge is in $T$, the only other information about $T$ that we need is a bound on $\|L_T\|\le2\Delta(T)$.
A forest drawn directly from a fractional split of $R$ may be a star.
So, before drawing $T$, we use randomized swap rounding~\cite{CVZ} to round $R$ to a random subgraph $F$ with about $\ell k$ edges and maximum degree $O(\ell)$, where $\ell=\Theta(\log n)$, and we draw $T$ from a fractional split of $F$.
This gives $\|L_T\|=O(\log n)$, and the matrix Freedman inequality keeps $L_R/q$ close to $B$ until $q=O(\log^2n)$.
We prove this in general form in Section~\ref{sec:extraction}.

For Theorem~\ref{thm:bounded}, we take $G=K_n$, $k=n-1$, $m=\lfloor n/2\rfloor$ and $B=L_{K_n}/(n/2)$.
For a set $X$ with at most $n/2$ vertices, closeness to $B$ gives $e_R(X)\le\frac35q(|X|-1)$, and, for a larger set $X$, it shows that many edges of $R$ meet $V\setminus X$; this gives every condition with room to spare.

For Theorem~\ref{thm:forests}, closeness to $B=L_G/m$ (suitably normalised by the degrees) only gives the conditions $e_R(X)\le q(|X|-1)$ if they hold in $G$ with some room, that is, if $e_G(X)\le(1-2\beta)m(|X|-1)$ for some $\beta>0$.
Following Xu~\cite{Xu}, we consider whether $G$ has an orientation in which every outdegree is at most $m-m^{1-\kappa}$, for a small constant $\kappa>0$.
If it does, then the room is there, with $\beta$ of order $m^{-\kappa}$, which is still enough for the argument above.
If it does not, then almost all edges of $G$ lie inside a set of about $k\ge1.9m$ vertices, and Xu handles this case with two further sampling arguments; we show in Appendix~\ref{app:xu} that they do not need properness.

\subsection{Exact decompositions}\label{sec:sketch-exact}
For Theorem~\ref{thm:gap}, we do not remove trees one at a time; instead, we build all of them together, and we use properness at every step.
\subsubsection{Small colour classes}\label{sec:sketch-linear}
We first prove Theorem~\ref{thm:gap} when every colour class has at most $2^{-49}n$ edges (Theorem~\ref{thm:linear} in Section~\ref{sec:recursion}).
Suppose that $n$ is even, so that we wish to decompose $E(K_n)$ into $n/2$ rainbow spanning trees.

\smallskip

\noindent\textbf{A recursive approach.}
A natural way to try to find such a decomposition is to divide and conquer.
Split the vertex set into two sets $A$ and $A'$ of even sizes $a$ and $a'$, and suppose that we have already decomposed $K[A]$ into $a/2$ rainbow spanning trees of $K[A]$, and $K[A']$ into $a'/2$ rainbow spanning trees of $K[A']$.
The remaining edges are those in $E(A,A')$, and their number is exactly right.
Indeed, split $A=A_1\cup A_2$ and $A'=A_1'\cup A_2'$ into halves, and let
\[
H=E(A_1,A_1')\cup E(A_2,A_2')\qquad\text{and}\qquad \barr{H}=E(A,A')\setminus H.
\]
Then each $v\in A'$ has exactly $a/2$ edges in $H$, that is, one for each tree on $A$, and each $u\in A$ has exactly $a'/2$ edges in $\barr{H}$, one for each tree on $A'$.
If, for each $v\in A'$, we give each of its edges in $H$ to a different tree on $A$, then each tree on $A$ gains exactly one pendant edge at each vertex of $A'$, and thus becomes a spanning tree of $K[A\cup A']$.
Doing the same with $\barr{H}$ for the trees on $A'$, we obtain $(a+a')/2$ edge-disjoint spanning trees of $K[A\cup A']$ which together use every edge of $K[A\cup A']$.
Repeating this, we can build a decomposition of $E(K_n)$ into spanning trees from decompositions of small complete subgraphs.
The only question is how to keep the trees rainbow.

\smallskip

\noindent\textbf{Adding the pendant edges at one step.}
Consider the trees $T_1,\ldots,T_q$ on $A$, where $q=a/2$, and, for each $v\in A'$, let $H(v)$ be the set of $q$ edges of $H$ at $v$.
We need, for each $v\in A'$, a bijection $\lambda_v:H(v)\to[q]$ indicating which tree receives which edge, so that each tree $T_i^+:=T_i\cup\{e:\lambda(e)=i\}$ is rainbow.
There are two ways in which a tree $T_i$ can fail to be rainbow after this step.
It can receive an edge whose colour it already contains, or it can receive two edges with the same colour at different vertices of $A'$.
(Two new edges at the same vertex always have different colours, as $\phi$ is proper.)
We choose each bijection $\lambda_v$ independently and uniformly at random.
Each of the failures above is then the event that one or two of these bijections take certain values.
Such events are close enough to independent that we can apply the lopsided local lemma, using a result of Lu and Sz\'ekely~\cite{LS} on random injections (see Lemma~\ref{lem:negdep}).
This works as long as no colour appears too often among the trees and the edges of $H$, and as long as the quantity
\begin{equation}\label{eq:sketch-r}
r:=\max_{v\in A',\,i\in[q]}\bigl|\phi(H(v))\cap \phi(T_i)\bigr|
\end{equation}
is at most a small fraction of $q$.
That is, no tree may already contain many of the colours of the edges it could receive at a vertex $v\in A'$.
If, for example, some tree $T_i$ contained all the colours of $H(v)$ except one, then $\lambda_v$ would have only one choice of which edge to give to $T_i$, and such forced choices for different trees could conflict.

\smallskip

\noindent\textbf{Controlling the colours of each tree throughout the recursion.}
We cannot simply assume that $r$ is small.
The trees on $A$ were themselves built at earlier steps of the recursion, and each step adds new colours to every tree.
As $\phi(H(v))\subseteq C(v)$, we have $r\leq \max_{v,i}\ell_{C(v)}(T_i)$, so it is enough to know that no tree contains many colours from any set $C(v)$.
We therefore include this in the statement we prove recursively.
At each step, we also ask that the random bijections do not give any tree many new edges with colour in $C(v)$, for any vertex $v$.
These are further events of the same kind as before, and the local lemma handles them as well.
For the deduction of the odd case of Theorem~\ref{thm:linear}, it is convenient to allow, in place of the sets $C(v)$, any family $\cZ$ of colour sets in which each $C(v)$ is contained in some member.
This leads to Lemma~\ref{lem:attach}. Roughly speaking, it says that, if no colour appears too often among the trees $T_i$ and the edges of $H$, and $r$ is small compared to $q$, then the bijections $\lambda_v$ can be chosen so that every new tree $T_i^+$ is rainbow.
Moreover, for each $Z\in \cZ$, fewer than $t$ of the new edges of $T_i^+$ have colour in $Z$, where $t$ depends on how many edges of $H$ have colour in $Z$.

As the number of colours in $Z$ used by a tree can increase by up to $t$ at each step, the recursion must have few levels, and, at each level, $t$ must be small compared to the size of the current vertex sets.
For the first requirement, we split each vertex set roughly in half, giving $O(\log n)$ levels.
For the second, we need that, inside each vertex set $W$ used in the recursion, no colour has many edges and no $Z\in \cZ$ has many edges with colour in $Z$.
To achieve this, we place the vertices into the recursion at random.

\smallskip

\noindent\textbf{Placing the vertices at random.}
We fix the shape of the recursion in advance, as a rooted binary tree $\cT$ whose nodes are subsets of $[n]$.
The root is $[n]$, the two children of a node partition it into two sets of about half its size, and the leaves of $\cT$ have size between $B$ and $2B$, for some $B=\Theta((\log n)^2)$.
We then choose a uniformly random bijection $\sigma:V(K_n)\to [n]$, and use the sets $\sigma^{-1}(W)$, $W\in V(\cT)$, as the vertex sets of the recursion.

To start the recursion, we need to decompose each $K[W]$, for $W$ a leaf of $\cT$, into rainbow spanning trees.
This is easy if $K[W]$ is itself rainbow: Walecki's construction from Section~\ref{sec:intro} decomposes $K[W]$ into $|W|/2$ Hamilton paths, which are then automatically rainbow.
For each $W$, we relabel the vertices of $W$ by a further random bijection, so that each of these Hamilton paths is a uniformly random Hamilton path, and hence contains few colours from each $Z\in\cZ$ (see Lemma~\ref{lem:ham}).

For a uniformly random $\sigma$, it is unlikely that $K[W]$ is rainbow for every leaf $W$ of $\cT$.
Indeed, for each pair of distinct edges with the same colour, we must avoid placing all four of their vertices inside the same leaf.
However, each of these events is again the event that $\sigma$ takes certain values on a few vertices, so the local lemma shows that, with positive probability, all of them are avoided.
Furthermore, the local lemma also bounds the probability of any other event of this form, conditioned on avoiding all of these events (see Lemma~\ref{lem:cond}).
This allows us to show that, conditioned on $K[W]$ being rainbow for every leaf $W$, with positive probability each set in the recursion contains few edges of each colour and few edges with colour in each $Z\in \cZ$.
Fixing such a bijection $\sigma$, we can then carry out the recursion, from the leaves of $\cT$ upwards, applying Lemma~\ref{lem:attach} at each step.

\subsubsection{Colour classes slightly below $n/2$}\label{sec:sketch-gap}
For Theorem~\ref{thm:gap}, we do not use the recursion of Section~\ref{sec:sketch-linear} on the whole of $K_n$. Instead, we add pendant edges only once, for many disjoint vertex sets in parallel.

\smallskip

\noindent\textbf{Trees inside random blocks.}
We partition the vertex set randomly into blocks $W_1,\ldots,W_k$, each of size about $X=(n\log n)^{1/3}$.
With high probability, every colour has only $O(\log n)$ edges inside each block, so we can apply Theorem~\ref{thm:linear} to decompose each $K[W_i]$ into $q_i=\lfloor |W_i|/2\rfloor$ rainbow spanning trees of $K[W_i]$.
There are $\sum_i q_i=\lfloor n/2\rfloor$ such trees in total, and it remains to turn each tree on $W_i$ into a spanning tree of $K_n$ by giving it one pendant edge at every vertex outside $W_i$.

\smallskip

\noindent\textbf{Sharing out the edges between blocks.}
First, we decide which edges go to which block.
For each $i$, every vertex $v\notin W_i$ must give exactly $q_i$ edges to the trees on $W_i$, one to each tree, and every edge between two blocks must go to at most one of the two blocks.
Moreover, as each tree on $W_i$ can use each colour at most once, no colour may be given to $W_i$ more than $q_i$ times.
To meet the degree requirements exactly, we pair up the vertices inside each block. For a pair $\{a,a'\}$ in $W_i$ and a pair $\{b,b'\}$ in $W_j$, the four edges between them form a copy of $K_{2,2}$, and we give one of its two perfect matchings to $W_i$ and the other to $W_j$.
Whatever choices we make, every vertex outside $W_i$ then gives exactly $q_i$ edges to $W_i$ (with a small modification for a block of odd size).
We choose a fractional assignment and round it using the Beck--Fiala theorem. This ensures that the number of edges of each colour given to $W_i$ is within $8$ of half of the number of edges of that colour available to $W_i$ (see Lemma~\ref{lem:quads}).

\smallskip

\noindent\textbf{Deciding which tree receives which edge.}
Let $E_i$ be the set of edges given to $W_i$.
We must decide which tree on $W_i$ receives each edge of $E_i$, so that each vertex outside $W_i$ gives one edge to each tree and each tree stays rainbow.
This is exactly a proper edge-colouring, with $q_i$ colours (one for each tree on $W_i$), of the bipartite graph $G_i$ with vertex classes $V(K_n)\setminus W_i$ and the set of colours of $\phi$, in which $v$ is joined to $\phi(vu)$ for each $vu\in E_i$.
A colour of $\phi$ which already appears on a tree inside $W_i$ must not be given to that tree, so this is a list edge-colouring problem.
We solve it in two stages using the theorem of Borodin, Kostochka and Woodall (see Lemma~\ref{lem:twostage}).

The gap in~\eqref{eq:gap} is used in this last step.
For the list colouring to exist, each colour may be given to $W_i$ at most $q_i-t$ times, where $t$ bounds, for each vertex $v\notin W_i$, the number of edges from $v$ to $W_i$ whose colour already appears inside $W_i$.
Remark~\ref{rem:twothirds} explains how this leads to the exponent $2/3$ in~\eqref{eq:gap}.
